\chapter{Những điều nên hiểu sớm}
\markboth{Những điều nên hiểu sớm}{Things to Understand Early}

\section{Điểm số không phải toàn bộ giá trị một con người.}
\begin{truyen}{Điểm số không phải toàn bộ giá trị một con người.}{A score is not a soul.}
\chuhoa{C}{ó học sinh xem một bài kiểm tra tệ như bằng chứng mình kém cỏi toàn diện. Nhiều năm sau nhìn lại, em mới thấy điểm số từng quan trọng, nhưng chưa bao giờ đủ để định nghĩa hết tính cách, lòng tốt, sức bền hay tương lai của mình.}
\textit{(Some students treat one bad test as proof that they are worthless. Years later, they see that scores mattered, but never defined their entire character, kindness, endurance, or future.)}
\end{truyen}
\begin{baihoc}
Đánh giá kết quả là cần, nhưng nhầm kết quả với con người là sai lầm rất sớm và rất nặng.
\textit{(Evaluating results is necessary, but confusing results with a person is an early and heavy mistake.)}
\end{baihoc}
\begin{ghinhoanh}
\textbf{Short English to remember:}

A score is not a soul.

Results matter, but they are not everything.
\end{ghinhoanh}
\begin{tuhocnhanh}
1. Em có đang dùng một kết quả để kết luận quá nhiều về bản thân không? \textit{(Are you using one result to conclude too much about yourself?)}

2. Ngoài điểm số, giá trị nào của em cần được nhìn lại? \textit{(Beyond scores, what value in you needs to be remembered?)}
\end{tuhocnhanh}
\ngancach

\section{Thông minh không thay thế được kiên trì.}
\begin{truyen}{Thông minh không thay thế được kiên trì.}{Talent without endurance fades.}
\chuhoa{N}{gày đi học, có người tiếp thu nhanh hơn hẳn bạn bè. Nhưng đi đường dài, người bền bỉ lại tiến xa hơn. Rất nhiều người chỉ hiểu điều này khi thấy mình từng có lợi thế mà vẫn dừng lại sớm.}
\textit{(In school, some understand much faster than others. But on long roads, the persistent often go further. Many people understand this only after watching early advantages fade.)}
\end{truyen}
\begin{baihoc}
Sự sắc bén giúp khởi đầu đẹp; sự kiên trì mới giữ được hành trình đẹp.
\textit{(Sharpness helps with a strong start; persistence sustains a strong journey.)}
\end{baihoc}
\begin{ghinhoanh}
\textbf{Short English to remember:}

Talent without endurance fades.

Stay longer than your mood.
\end{ghinhoanh}
\begin{tuhocnhanh}
1. Em đang dựa nhiều vào năng khiếu hay vào sự bền bỉ? \textit{(Are you relying more on talent or on persistence?)}

2. Việc gì em cần theo đủ lâu mới thấy kết quả? \textit{(What do you need to stay with long enough to see results?)}
\end{tuhocnhanh}
\ngancach

\section{Thầy cô cũng từng sai, và người lớn cũng đang học dở dang.}
\begin{truyen}{Thầy cô cũng từng sai, và người lớn cũng đang học dở dang.}{Adults are unfinished too.}
\chuhoa{N}{hiều học sinh tưởng lớn lên là hết lúng túng. Nhưng càng trưởng thành, con người càng thấy mình còn thiếu nhiều điều. Hiểu được điều đó sớm giúp ta bớt thần tượng hóa và cũng bớt cay nghiệt với nhau hơn.}
\textit{(Many students think adulthood means certainty. But the older people grow, the more they see what they still lack. Understanding this early makes us less naïve and also less cruel to one another.)}
\end{truyen}
\begin{baihoc}
Biết người lớn cũng bất toàn giúp ta trưởng thành với nhiều cảm thông hơn.
\textit{(Knowing that adults are also imperfect helps us mature with more compassion.)}
\end{baihoc}
\begin{ghinhoanh}
\textbf{Short English to remember:}

Adults are unfinished too.

Growing up does not end learning.
\end{ghinhoanh}
\begin{tuhocnhanh}
1. Em có đang đòi người lớn quanh mình phải hoàn hảo không? \textit{(Are you expecting the adults around you to be perfect?)}

2. Nhìn người khác bất toàn giúp em sống cảm thông hơn ra sao? \textit{(How can seeing others as imperfect help you live more compassionately?)}
\end{tuhocnhanh}
\ngancach

\section{Không phải ai chăm chỉ cũng thành công, nhưng lười thì gần như chắc chắn không.}
\begin{truyen}{Không phải ai chăm chỉ cũng thành công, nhưng lười thì gần như chắc chắn không.}{Effort is not a guarantee, but a baseline.}
\chuhoa{C}{uộc đời không công bằng tới mức ai chăm chỉ cũng thắng lớn. Nhưng nó đủ rõ ở một điểm: lười biếng hiếm khi đưa ai đến nơi đáng đến. Hiểu sớm điều này giúp ta bớt ảo tưởng và bớt buông xuôi.}
\textit{(Life is not so fair that every hard worker wins big. But it is clear about one thing: laziness rarely takes anyone somewhere worth reaching. Learning this early reduces illusion and excuses.)}
\end{truyen}
\begin{baihoc}
Nỗ lực không hứa phần thưởng tuyệt đối, nhưng thiếu nỗ lực thì gần như tự loại mình từ đầu.
\textit{(Effort does not promise absolute reward, but lack of effort almost eliminates you from the start.)}
\end{baihoc}
\begin{ghinhoanh}
\textbf{Short English to remember:}

Effort is not a guarantee, but a baseline.

Do the work before judging the world.
\end{ghinhoanh}
\begin{tuhocnhanh}
1. Em có đang thất vọng về cuộc đời khi bản thân còn chưa làm đủ không? \textit{(Are you disappointed with life before you have really done enough?)}

2. Nỗ lực tối thiểu nào em còn đang thiếu? \textit{(What minimum effort are you still missing?)}
\end{tuhocnhanh}
\ngancach

\section{Học là cho mình, không phải chỉ để làm vừa lòng người khác.}
\begin{truyen}{Học là cho mình, không phải chỉ để làm vừa lòng người khác.}{Borrowed motivation does not last.}
\chuhoa{C}{ó người học vì sợ bị mắng, vì muốn được khen, vì không muốn thua bạn. Những động lực ấy giúp đi một đoạn ngắn. Đi lâu, chỉ người hiểu rằng học là để mở đời mình mới đủ bền để tiếp tục.}
\textit{(Some people study to avoid scolding, to gain praise, or to beat friends. Those motives can carry them for a short while. On longer roads, only those who know learning expands their own life can continue steadily.)}
\end{truyen}
\begin{baihoc}
Khi động lực nằm ngoài mình hết, việc học sẽ rất dễ gãy khi hoàn cảnh đổi.
\textit{(When motivation lies entirely outside us, learning breaks easily when circumstances change.)}
\end{baihoc}
\begin{ghinhoanh}
\textbf{Short English to remember:}

Borrowed motivation does not last.

Learn for your own life.
\end{ghinhoanh}
\begin{tuhocnhanh}
1. Em đang học vì điều gì sâu nhất? \textit{(What is the deepest reason you are learning?)}

2. Nếu không ai chấm điểm, em còn muốn giỏi hơn ở điều gì? \textit{(If no one graded you, what would you still want to become better at?)}
\end{tuhocnhanh}
\ngancach